\section{Introducción al Dossier de Firmware}

Este dossier presenta la arquitectura y especificación del firmware y software embebido del collar de geocercamiento DR-VF-01. Se detallan la secuencia de arranque, la máquina de estados finitos (FSM), las estrategias de mitigación de latencia en tiempo real y la definición del firmware mínimo viable (MVP).

El firmware es un único código base sobre \textbf{Zephyr RTOS} para dos placas: la prueba de concepto en la \textbf{Nordic Thingy:91 X} (nRF9151, nRF Connect SDK, LTE-M/NB-IoT/GNSS) y el producto final en la \textbf{PCB con ESP32-S3-WROOM-1U-N16R8} (LoRa SX1262, GNSS externo MAX-M10S, IMU LSM6DSOX). Se comparten \texttt{geofence.c}, la FSM y el registro de eventos (\emph{logging}); cada placa aporta su \emph{overlay} de devicetree, su \texttt{prj.conf} y los adaptadores de GNSS, comunicaciones e IMU.

\section{Secuencia de Arranque (\textit{Power-on Sequence})}

\subsection{Fases de Inicialización}

\begin{enumerate}[leftmargin=*]
    \item \textbf{Power-On Reset (POR)}: Verificación de estabilización de riel de alimentación ($3.3\text{ V}$) y reinicio del SoC.
    \item \textbf{Inicialización de Periféricos}: Configuración de buses de comunicación ($\text{I}^2\text{C}$, SPI) y líneas GPIO de interrupción.
    \item \textbf{Inicialización de Sensores}: Autochequeo (BIST) y configuración del transceptor LoRa SX1262, IMU LSM6DSOX y GNSS MAX-M10S.
    \item \textbf{Carga de Parámetros de Geocerca}: Lectura del polígono activo y llaves criptográficas desde memoria Flash NVS.
    \item \textbf{Arranque del RTOS (Zephyr)}: Creación de tareas prioritarias y mecanismos de sincronización (\textit{semáforos y colas}).
\end{enumerate}

\subsection{Presupuesto de Tiempo de Arranque}

El tiempo total de arranque desde el evento de energización hasta el ingreso al bucle operativo principal es de $5.0\text{ s}$:

\begin{itemize}[leftmargin=*]
    \item \textbf{POR y calibración básica de reloj}: $1.2\text{ s}$
    \item \textbf{Configuración de buses periféricos}: $1.8\text{ s}$
    \item \textbf{Inicialización y BIST de sensores}: $1.5\text{ s}$
    \item \textbf{Carga NVS e inicio de RTOS}: $0.5\text{ s}$
\end{itemize}

\section{Máquina de Estados Finitos (FSM)}

\subsection{Estados Principales del Sistema}

\begin{figure}[H]
\centering
\resizebox{\textwidth}{!}{%
\begin{tikzpicture}[
    node distance=2.2cm and 3.6cm,
    auto,
    >=stealth,
    state/.style={
        draw=bordergrey,
        thick,
        rectangle,
        rounded corners=6pt,
        fill=categorybg,
        text width=2.8cm,
        minimum height=1.1cm,
        align=center,
        font=\small\bfseries
    },
    critical/.style={
        state,
        fill=phase2orange,
        draw=unalred
    },
    errorstate/.style={
        state,
        fill=categorybgsub,
        draw=unalnavy
    },
    line/.style={
        draw=unalnavy,
        thick,
        ->
    }
]

    % --- NODOS PRINCIPALES ---
    \node [state] (init) {\texttt{INIT}};
    \node [state, right=of init] (sleep) {\texttt{SLEEP}};
    \node [state, right=of sleep] (check) {\texttt{GEOFENCE\_CHECK}};
    \node [state, right=of check] (gnss) {\texttt{GNSS\_UPDATE}};

    % --- NODOS DE ALERTA Y ERROR (FILA INFERIOR) ---
    \node [errorstate, below=of init] (error) {\texttt{ERROR}};
    \node [critical, below=of sleep] (disuasion) {\texttt{DISUASION}};
    \node [critical, below=of check] (warning) {\texttt{WARNING}};

    % --- TRANSICIONES ARRANQUE Y CONTROL ---
    \path [line] (init) -- node[above, font=\scriptsize, align=center] {POR / BIST OK} (sleep);
    \path [line] (init) -- node[left, font=\scriptsize, align=center] {BIST Fail} (error);

    % --- BUCLE DE BAJO CONSUMO (SLEEP <-> CHECK) ---
    \path [line] (sleep) edge[bend left=22] node[above, font=\scriptsize, align=center] {RTC / IMU Wakeup} (check);
    \path [line] (check) edge[bend left=22] node[below, font=\scriptsize, align=center] {In-Fence / Sleep} (sleep);

    % --- BUCLE GNSS UPDATE (HOT START) ---
    \path [line] (check) edge[bend left=22] node[above, font=\scriptsize, align=center] {Stale Fix} (gnss);
    \path [line] (gnss) edge[bend left=22] node[below, font=\scriptsize, align=center] {Fix Valid} (check);

    % --- CICLO DE ADVERTENCIA Y DISUASIÓN ---
    \path [line] (check) -- node[right, font=\scriptsize, align=center] {Near Border} (warning);
    \path [line] (warning) -- node[below, font=\scriptsize, align=center] {Breach\\[-1pt]{\tiny [Dist $>$ Limit]}} (disuasion);
    \path [line] (disuasion) -- node[left, font=\scriptsize, align=center] {Secured /\\[-1pt]Timer Expired} (sleep);

    % --- RUTA DEFENSIVA A ERROR (Rodea INIT por el exterior izquierdo) ---
    \draw [line] (check.north) -- ++(0, 1.5) -| ([xshift=-1.0cm]error.west) node[pos=0.3, above, font=\scriptsize, align=center] {WDT Expired / Hardware Fault} -- (error.west);

\end{tikzpicture}%
}
\caption{Máquina de Estados Finitos Defensiva (FSM) del collar DR-VF-01 con soporte para ráfagas GNSS y recubrimiento WDT (según el modelo de Elecia White)}
\label{fig:fsm_diagram_complete}
\end{figure}

\subsection{Descripción Formal de los Estados}

\begin{description}[leftmargin=*, style=nextline, font=\bfseries\ttfamily]
    \item[INIT] Secuencia de arranque, verificación de integridades en Flash/RAM (BIST) y preparación del Kernel RTOS.
    \item[SLEEP] Modo de bajo consumo activo con dominios de energía secundarios deshabilitados ($I_q < 20\,\mu\text{A}$).
    \item[GEOFENCE\_CHECK] Ejecución en tiempo real del algoritmo Ray-Casting optimizado para evaluación del perímetro.
    \item[GNSS\_UPDATE] Ráfaga corta de adquisición GNSS en modo \textit{Hot Start} ante pérdida temporizada de precisión HDOP.
    \item[WARNING] Emisión de tono acústico progresivo ante aproximación al límite ($< 5\,\text{m}$).
    \item[DISUASION] Activación de la secuencia de impulsos disuasivos corregidos tras rebasar la geocerca.
    \item[ERROR] Modo defensivo de aislamiento por desbordamiento de temporizador WDT o falla en buses $\text{I}^2\text{C}$/SPI.
\end{description}

\subsection{Transición Crítica y Latencia}

La transición más crítica del sistema ocurre entre los estados \texttt{SLEEP} y \texttt{GEOFENCE\_CHECK}. Tras una interrupción externa generada por la IMU, la FSM debe reanudar la ejecución y evaluar la posición en menos de $50\,\text{ms}$.

\section{Análisis de Desafíos en Tiempo Real}

\subsection{Restricción Temporal: Algoritmo Point-in-Polygon ($< 50\,\text{ms}$)}

\subsubsection{Implementación Optimizada del Algoritmo Ray-Casting}

El algoritmo PiP se implementó en C puro optimizando las operaciones de punto flotante para aprovechar la unidad FPU hardware del ESP32-S3:

\begin{verbatim}
bool pip_optimized(const point_t* p, const polygon_t* poly) {
    int crosses = 0;
    const int vertex_count = poly->num_vertices;
    
    for (int i = 0; i < vertex_count; i++) {
        const point_t* v1 = &poly->vertices[i];
        const point_t* v2 = &poly->vertices[(i + 1) % vertex_count];
        
        // Condición de descarte rápido por límites verticales
        if ((v1->y > p->y) == (v2->y > p->y)) continue;
        
        // Cálculo de intersección optimizado para FPU
        const float vt = (p->y - v1->y) / (v2->y - v1->y);
        if (p->x < v1->x + vt * (v2->x - v1->x)) {
            crosses++;
        }
    }
    
    return (crosses & 1); // Verdadero si el número de cruces es impar
}
\end{verbatim}

\subsubsection{Gestión de Memoria y PSRAM}

Para evitar cuellos de botella por latencia de acceso a Flash, la estructura de datos del polígono activo (hasta 64 vértices) se aloja en PSRAM bus octal de alta velocidad.

\subsection{Interrupción de Despertar (\textit{GPIO Wakeup}) desde IMU}

\subsubsection{Configuración del Registro del LSM6DSOX}

\begin{verbatim}
int lsm6dsox_configure_wakeup(void) {
    // Configurar umbral de detección de movimiento (0.5g)
    lsm6dso_write_reg(LSM6DSO_WAKE_UP_THS, 0x20);
    
    // Configurar duración mínima (16 muestras consecutivas)
    lsm6dso_write_reg(LSM6DSO_WAKE_UP_DUR, 0x10);
    
    // Habilitar la ruta de interrupción en el pin INT1
    lsm6dso_write_reg(LSM6DSO_INT1_CTRL, 0x20);
    
    // Ajustar tasa ODR en modo bajo consumo (26 Hz)
    lsm6dso_write_reg(LSM6DSO_CTRL1_XL, 0x10);
    
    return 0;
}
\end{verbatim}

\section{Firmware Mínimo Viable (MVP)}

\subsection{Módulos Core del MVP}

\begin{enumerate}[leftmargin=*]
    \item \textbf{Gestor de Movimiento}: Driver IMU LSM6DSOX con detección de umbral de aceleración.
    \item \textbf{Motor de Geocercamiento}: Algoritmo Ray-Casting optimizado con almacenamiento en RAM.
    \item \textbf{Pila LoRaWAN / LoRa P2P}: Transmisión de eventos prioritarios mediante SX1262.
    \item \textbf{Interfaz de Configuración}: Punto de acceso SoftAP con servidor HTTP embebido.
    \item \textbf{Administrador de Energía}: Gestión de modos \textit{Deep Sleep} e interrupciones del sistema.
\end{enumerate}

\subsection{Definición de Hilos en RTOS}

\begin{verbatim}
K_THREAD_DEFINE(imu_monitor, STACK_SIZE, imu_monitor_thread, 
                NULL, NULL, NULL, IMU_PRIORITY, 0, 0);

K_THREAD_DEFINE(geofence_check, STACK_SIZE, geofence_thread,
                NULL, NULL, NULL, GEOFENCE_PRIORITY, 0, 0);

K_THREAD_DEFINE(lora_handler, STACK_SIZE, lora_thread,
                NULL, NULL, NULL, LORA_PRIORITY, 0, 0);
\end{verbatim}

\section{Conclusiones}

La arquitectura de firmware propuesta para el collar DR-VF-01 combina un esquema determinista multitarea con una gestión rigurosa de bajo consumo. La optimización del algoritmo Point-in-Polygon y el uso eficiente de interrupciones por hardware garantizan tiempos de respuesta inferiores a $50\,\text{ms}$, asegurando el funcionamiento correcto de las acciones de disuasión antes de que el ganado traspase el perímetro configurado.